\documentclass[10pt]{article}

\usepackage[utf8]{inputenc}

\usepackage[T1]{fontenc}

\usepackage{amsmath}

\usepackage{amsfonts}

\usepackage{amssymb}

\usepackage[version=4]{mhchem}

\usepackage{stmaryrd}

\usepackage{graphicx}

\usepackage[export]{adjustbox}

\graphicspath{ {./images/} }

\usepackage{bbold}



\begin{document}

все алгебраич. подмногообразия $X$ размерности $r$ и степени $d$ шроективного пространства $P^{n}$.



В произведении $X \times\left(\check{P}^{n}\right)^{\boldsymbol{r}+1}$, где $\check{\boldsymbol{P}}^{n}$-двойственное к $P^{n}$ проективное пространство, параметризующее гиперплоскости $u \subset P^{n}$, рассматривается подмногообразие



\[

\Gamma=\left\{x, u^{(0)}, \ldots, u^{(r)} \mid x \in u^{(i)} \text { при } i=0, \ldots, r\right\} .

\]



Его образ $p_{2}(\Gamma) \subset(\check{P} n)^{r+1}$ при проекции на второй сомножитель есть гиперповерхность в $\left(\breve{\boldsymbol{P}}^{n}\right)^{r+1}$, к-рая задается формой $F_{X}$ от $r+1$ системы и по $n+1$ переменным, однородной степени $d$ по каждой системе переменных. Форма $F_{X}$ наз. а с с о ц и и р о в а в в о й фо р м о й (или форм о й К эли) многообразия $X$; она полностью определяет подмногообразие $X$. Эта форма была введена Б. Л. Ван дер Варденом и В. Чжоу [1]. Коэфффициенты формы $F_{X}$ определены с точностью до постоянного множителя и наз. к о о р д и натам и Ч ж о у многообразия $X$.



Координаты Чжоу многообразия $X$ определяют точку $c(X) \in P^{v}$, где $v$-нек-рая функция от $n, r, d$. Точки $c(X) \in P^{v}$, соответствующие всем неприводимым подмногообразиям $X \subset P^{n}$ размерности $r$ и степени $d$, заполняют в $P^{v}$ квазипроективное подмногообразие $C_{n, r, d}$, называемое м в о г о об р а з и м Ч жо у. Если рассматривать не только неприводимые подмногообразия, но и положительные алгебраич. циклы (т. е. формальные линейные комбинации многообразий с целыми положительными коэффициентами) размерности $r$ и степени $d$ в $P^{n}$, то получается замкнутое подмногообразие $\bar{C}_{n, r, d} \subset P^{v}$, к-рое также наз. многообразием Чжоу. Ч. м. является базой универсального алгебраич. семейства $\mathfrak{X} \stackrel{\pi}{\longrightarrow} \bar{C}_{n, r, d}$, где $\mathfrak{X} \subset \bar{C}_{n, r, d} \times P^{n}, \pi$ индуцировано проекцией и слой $\pi^{-1}(c)$ в точке $c(X) \in \bar{C}_{n, r}^{\prime}, d$ совпадает с циклом $X$. Простейшими примерами Ч. М. являются многообразия $C_{3,1, d}$-кривых степени $d$ в $P^{3}$. Так, $C_{3,1,1}=\bar{C}_{3,1,1}$ - неприводимое многообразие размерности 4, изоморфное квадрике Плюккера в $P^{\bar{s}}$; $\bar{C}_{3,1,2}=C^{(1)} \cup C^{(2)}$ состоит из двух компонент размерности 8, где $C^{(1)}$ соответствует плоским кривым 2-го порядка, a $C^{(2)}$-парам прямых; $\bar{C}_{3,1,3}$ состоит из четырех компонент размерности 12 , к-рые соответствуют тройкам прямых, кривым, состоящим из прямой и плоской квадрики, плоским кубикам, неплоским кривым степени 3. Во всех этих случаях многообразия $C_{3,1, d}$ рациональны. Однако из нерациональности схемы модулей кривых достаточно большого рода следует, что при достаточно больших $d$ многообразия $C_{3,1, d}$ нерациональны (см. [2]).



Если $V \subset P^{n}$-алгебраич. подмногообразие, то циклы $Z \subset P^{n}$ размерности $r$ и степени $d$, лежащие в $V$, образуют алгебраич. подмногообразие $\bar{C}_{r, d}(V) \subset \bar{C}_{n, r} d$. Этот результат позволяет ввести нек-рую алгеброгеометрич. структуру на множестве положительных $r$-мерных циклов $Z_{r}^{+}(V)=U_{d}>{ }_{0} \bar{C}_{r, d}(V)$ многообразня $V$ (см. [1]).



О других подходах к проблеме классификации многообразий см. Гильберта схема, Модулей проблема.



Jum.: [1] V a n der W a e r d e n B. L. $\mathrm{C}$ h o w W.-L., "Math. Ann.», 1937, Bd 113, S. $692-704 ;$ [2] H a r r i s J.,

M u m o r d D., "Invent. Math». 1982, v. 67, p. $23-86:$ [3] Х о д ж В., П и д о Д., Методы алгебраической геометрии, пер. с англ., т. 2, М., 1954; [4] IШ а ф а р е в ич И. Р., Основы алгеб-

раической геометрии, М., 1972.

Вал. С. Куликов. раической геометрии, М., 1972.



ЧЖОУ ТЕОРЕМА: любое аналитич. подмножество комплексного проективного пространства является алігебраическим многообразием. Теорема доказана В. Чжо [1].



Лum.: [1] C h o w W.-L., "Amer. J. Math.", 1949, v. 71, p. $893-914 ;$ [2] Г р и фофи т с Ф., Х а р р и с Д к., Принципы



\begin{center}

\includegraphics[max width=\textwidth]{2023_07_09_5608f5b3596dbefdce07g-1}

\end{center}



ЧЖӘНЯ КЛАСС-характеристический класс, определенный для комплексных векторных расслоений. Ч. к. комплексного векторного расслоения $\xi$ с базой $B$ обозначается $c_{i}(\xi) \in H^{2 i}(B)$ и ошределен для всех натуральных индексов $i$. П о л н м Ч. к. наз. неоднородный характеристич. класс $c=1+c_{1}+c_{2}+\ldots$, а по ли н о м о м Ч ж э н я-выражение $c_{t}=1+c_{1} t+$ $+c_{2} t^{2}+\ldots$, где $t$ - формальная переменная. Ч. к. введены в [1].



Характеристич. классы, определенные для $n$-мерных комплексных векторных расслоений, со значениями в целотисленных когомологиях, естественно отождествляются с әлементами кольца $H^{* *}\left(B U_{n}\right)$. В этом смысле Ч. к. $c_{i}$ можно считать элементами группы $H^{2 i}\left(B U_{n}\right)$, полный Ч. к.- неоднородным әлементом кольца $H^{* *}\left(B U_{n}\right)$, а полином Чжэня-элементом кольца формальных степенных рядов $H^{* *}\left(B U_{n}\right)[[t]]$



\begin{enumerate}

  \setcounter{enumi}{3}

  \item к. обладают следующими свойствами, к-рые однознатно их определяют. 1) Для векторных расслоений $\xi, \eta$ с общей базой $B c(\xi \oplus \eta)=c(\xi) c(\eta)$, другими словами, $\left.c_{k}(\xi \bigoplus \eta)=\sum_{i} c_{i}(\xi) c_{k-i}(\eta), c_{0}=1.2\right)$ Для одномерного универсального расслоения $\varkappa_{1}$ над $\mathbb{C} P \infty$ имеет место равенство $c\left(x_{1}\right)=1+u$, где $\left.u \in H^{2}(\mathbb{C P})^{\infty}\right)$-ориентация расслоения $x_{1}\left(\mathbb{C} P_{\infty}-\right.$ Tома пространство расслоения $x_{1}$, к-рое, будучи комплексным, имеет однознатно определенную ориентацию $u)$.

\end{enumerate}



Следствия свойств $1-2): c_{i}(\xi)==0$ при $i>\operatorname{dim} \xi ;$ $c(\xi)=c(\xi \oplus \theta)$, где $\theta$-тривиальное расслоение. Последнее обстоятельство позволяет определить Ч. к. как элементы колыца $H^{* *}(B U)$.



Если $\omega=\left\{i_{1}, \ldots, i_{k}\right\}$ - набор целых неотрицательных чисел, то через $c_{\omega}$ обозначается характеристич. класс $c_{i_{1}} c_{i_{2}} \ldots c_{i_{k}} \in H^{2 n}(B U)$, где $n=i_{1}+\ldots+i_{k}$



При естественном мономорфизме $H^{* *}\left(B U_{n}\right) \rightarrow$ $\longrightarrow H^{* *}\left(B T_{n}\right)=\mathbb{Z}\left[\left[x_{1}, \ldots, x_{n}\right]\right]$, индуцированном отображением $\stackrel{n}{B}_{n}=\mathbb{C} P^{\infty} \times \ldots \times \mathbb{C}^{\infty} \rightarrow \cdots B U_{n}$, Ч. к. переходят в әлементарные симметрич. функщии, а полный Ч. к. переходит в многочлен $\prod_{i=1}^{n}\left(1+x_{i}\right)$. Образом кольца $H^{* *}\left(B U_{n}\right)$ в $H^{* *}\left(B T_{n}\right)=\mathbb{Z}\left\lceil\left[x_{1}, \ldots, x_{n}\right]\right]$ является годкольцо всех симметрических формалыных степенных рядов. Каждый симметрическшй формальный стеленной ряд от опразуюших Ву $x_{1}, \ldots, x_{n}$ определяет характеристич. класс, к-рый может быть выражен через Ч. к. Напр., ряд $\prod_{i=1}^{n} \frac{x_{i}}{1-\mathrm{e}^{x_{i}}}$ определяет характеристич. класс с рацгональными коюффициентами, наз. кла с со м То дла и обозначается $T \in$ $\in H^{* *}\left(B U_{n} ; \mathbb{Q}\right)$.



Гугть $\omega=\left\{i_{1}, \ldots, i_{k}\right\}$ - набор пелых неотрицательных чисел. Через $\dot{S}_{\omega}\left(c_{1}, \ldots, c_{n}\right)$ обозначается характеристич. класс, определяемый наименышим симметрич. многочленом от переменных $x_{1}, \ldots, x_{n}$, где $n \geqslant i_{1}+\ldots+$ $+i_{k}$, содержащим одночлен $x_{1}^{i_{1}} \ldots x_{k}^{i_{k}}$.



Пусть $h^{*}$-ориентированная мультипликативная теория когомологий. Ч. к. $\sigma_{i}$ со значениями в теории $h^{*}$ обладают такими же, что п обычные Ч. к., свойствами: $\sigma(\xi \oplus \eta)=\sigma(\xi) \sigma(\eta), \quad \sigma=1+\sigma_{1}+\sigma_{2}+\ldots ; \sigma\left(x_{1}\right)=1+$ $+u \in h^{*}\left(\mathbb{C} P^{\infty}\right)$, где $u \in h^{2}\left(\mathbb{C} P^{\infty}\right)$-ориентация расслоения $x_{1}$, к-рые их однозначно определяют. Как и для обычных Ч. к., уІотребляются обозначения $\sigma_{\omega}=$ $=\sigma_{i_{1}} \ldots \sigma_{i_{k}}$ и $S_{\omega}\left(\sigma_{1}, \ldots, \sigma_{n}\right)$. Если $\xi, \eta$ - комплексные векторные расслоения, то



\[

=\begin{gathered}

S_{\omega}\left(\sigma_{1}, \ldots, \sigma_{n}\right)(\xi \oplus \eta)= \\

=\sum_{\omega^{\prime} U \omega^{\prime \prime}=\omega^{\prime}} S_{\omega^{\prime}}\left(\sigma_{1}, \ldots, \sigma_{n}\right)(\xi) S_{\omega^{\prime \prime}}\left(\sigma_{1}, \ldots, \sigma_{n}\right)(\eta),

\end{gathered}

\]





\end{document}